
\uline{entier} $n$ \uline{tel que pour tout} $k\ge1$, \uline{il existe une section} $s\in\Gamma(X,\underline{E}(kn))$ \uline{qui soit une base de} $\underline{E}(kn)_{x_i}$ \uline{pour tout} $i$.

On peut trouver un entier $n_0$, et une section $f\in\Gamma(X,\underline{L}^{\otimes n_0})$ telle que pour tout $i$, $x_i\in X_f$, et que $X_f$ soit un ouvert affine (EGA II 4.5.4). Sur $X_f$, on peut trouver une section $t\in\Gamma(X_f,\underline{E})$ satisfaisant les conditions du lemme 6.7; de plus, il existe un entier $n'_0$ tel que la section $f^{\otimes n'_0}\otimes t\in\Gamma(X_f,\underline{E}(n_0n'_0))$ se prolonge à $X$ tout entier. Posons $n=n_0n'_0$; alors $f^{\otimes n'_0}\in\Gamma(X,\underline{L}^n)$; en faisant le produit tensoriel de la section de $\underline{E}(n)$ obtenue plus haut et de $f^{\otimes kn'_0}$, on obtient donc une section de $\underline{E}((k+1)n)$ pour tout $k\ge0$. Comme pour tout $p$, $X_{f^{\otimes p}}=X_f$, et que la section $t$ est une base de $\underline{E}_{x_i}$ pour tout $i$, les sections de $\underline{E}(kn)$ ainsi construites pour tout $k\ge1$ forment une base de $\underline{E}(kn)_{x_i}$ pour tout $i$.

\textbf{Théorème 6.9.} \uline{Soit} $X$ \uline{un schéma noethérien de dimension} $d$, \uline{ayant un faisceau ample} $\underline{L}$. \uline{Alors} $\Fil^{d+1}K^{\bullet}(X)=0$, \uline{et à fortiori} $(\Fil^1K^{\bullet}(X))^{d+1}=0$.

(Cette dernière relation est due à J.P. Serre, dans le cas général envisagé ici).

Nous allons d'abord montrer que si $\underline{E},\underline{E}'$ sont deux faisceaux inversibles sur $X$, de classes $E,E'$, alors $E-E'\in K^{\bullet}(X)_1$.

Soient $Y$ un fermé de $X$, $Y_1,\ldots,Y_k$ ses composantes irréductibles, et pour tout $i$, soit $x_i$ un point de $Y_i-\bigcup_{j\ne i}Y_j$. D'après le lemme 6.8, il existe des entiers $n,n'$, tels que, quels que soient les entiers $k,k'\ge1$, on puisse trouver des homomorphismes de $\underline{O}_X$ dans $\underline{E}(kn)$, $\underline{E}'(k'n')$ qui soient des isomorphismes aux points $x_i$; on obtient donc des morphismes:
\[
  \underline{L}^{\otimes -kn}\xrightarrow{s_k} \underline{E},
  \qquad
  \underline{L}^{\otimes -k'n'}\xrightarrow{s'_{k'}} \underline{E}',
\]
